\documentclass[12pt]{article}

% === Packages ===
\usepackage[margin=1in]{geometry}
\usepackage{amsmath, amssymb}
\usepackage{setspace}
\usepackage{titlesec}
\usepackage{enumitem}
\usepackage{helvet}
\renewcommand{\familydefault}{\sfdefault} % modern sans-serif look

% === Formatting tweaks ===
\setstretch{1.15}
\setlist[enumerate]{itemsep=0.4em, topsep=0.2em}
\titleformat{\section}{\Large\bfseries\sffamily}{\thesection.}{1em}{}
\titleformat{\subsection}{\large\bfseries\sffamily}{\thesubsection}{1em}{}

% === Title ===
\title{\vspace{-2em}\textbf{Applied Mathematics Oral Exam Questions and Model Answers}\\
\large From Euler to Navier--Stokes}
\author{Prepared for Graduate Oral Examination}
\date{}

\begin{document}
\section*{Oral Exam Questions and Model Answers Based on the Solid Mechanics Tutorial}

%%%%%%%%%%%%%%%%%%%%%%%%%%%%%%%%%%%%%%%%%%%%%%%%%%%%%%%
\subsection*{1. What is stress, and how is it defined in 1D and 3D?}

\textbf{Answer.}
In 1D, stress is defined as force per area:
\[
\sigma = \frac{F}{A}.
\]


In 3D, stress becomes a \textbf{second-order tensor} whose components 
$\sigma_{ij}$ represent the $i$-direction force acting on a surface with 
normal in the $j$-direction. Stress in 3D is written as
\[
\boldsymbol{\sigma} =
\begin{bmatrix}
\sigma_{11} & \sigma_{12} & \sigma_{13} \\
\sigma_{21} & \sigma_{22} & \sigma_{23} \\
\sigma_{31} & \sigma_{32} & \sigma_{33}
\end{bmatrix}.
\]
This tensorial formulation is necessary because real materials experience 
forces in multiple directions simultaneously.


%%%%%%%%%%%%%%%%%%%%%%%%%%%%%%%%%%%%%%%%%%%%%%%%%%%%%%%
\subsection*{2. What is strain, and what does it measure?}

\textbf{Answer.}
Strain measures deformation:
\[
\epsilon = \frac{\Delta L}{L_0}
\]


In multiple dimensions, strain is also a tensor reflecting how material
line elements deform along and across coordinate directions.  
The full strain tensor is used to capture shear, normal extension,
and multidimensional changes in shape.


%%%%%%%%%%%%%%%%%%%%%%%%%%%%%%%%%%%%%%%%%%%%%%%%%%%%%%%
\subsection*{3. Explain Hooke’s law.}

\textbf{Answer.}
In 1D, Hooke’s law is
\[
\sigma = E \epsilon,
\]
where $E$ is the Young’s modulus (p.~9–10).


%%%%%%%%%%%%%%%%%%%%%%%%%%%%%%%%%%%%%%%%%%%%%%%%%%%%%%%
\subsection*{4. What types of stress are shown in the tutorial, and why do they matter in biomechanics?}

\textbf{Answer.}


In biomechanics:
\begin{itemize}
    \item \textbf{Tensile}: ligaments and tendons (stretching).
    \item \textbf{Compressive}: bone and cartilage (weight-bearing).
    \item \textbf{Shear}: tissue sliding, muscle layers, arterial walls.
\end{itemize}
IBFE simulations often must capture all three simultaneously, especially
for biological soft tissues which experience shear and tension together.


%%%%%%%%%%%%%%%%%%%%%%%%%%%%%%%%%%%%%%%%%%%%%%%%%%%%%%%
\subsection*{5. Define traction and explain how it differs from pressure.}

\textbf{Answer.}
Traction is the force per unit area acting on a surface with a specified

It may include normal and tangential components.

Pressure is always normal to the surface, whereas traction may have 
shear components.  
In IBFE, the traction from the fluid onto the solid determines the
boundary loading on the Lagrangian solid mesh.


%%%%%%%%%%%%%%%%%%%%%%%%%%%%%%%%%%%%%%%%%%%%%%%%%%%%%%%
\subsection*{6. What are the three main types of mechanical loading shown in the tutorial?}

\textbf{Answer.}


\begin{itemize}
    \item \textbf{Axial}: produces uniform tensile or compressive stress.
    \item \textbf{Bending}: produces tension on one side and compression on the other.
    \item \textbf{Torsion}: produces shear stress.
\end{itemize}

In IBFE simulations, immersed structures (e.g. beams, muscles, tendons)
experience all of these depending on fluid forces and constraints.


%%%%%%%%%%%%%%%%%%%%%%%%%%%%%%%%%%%%%%%%%%%%%%%%%%%%%%%
\subsection*{7. Describe Young’s modulus, Poisson’s ratio, and yield strength.}


\begin{itemize}
    \item \textbf{Young's modulus $E$}: stiffness; slope of linear portion of the stress–strain curve.
    \item \textbf{Poisson ratio $\nu$}: ratio of lateral to axial strain,
    \[
    \nu = -\epsilon_{\text{lateral}}/\epsilon_{\text{axial}}.
    \]
    \item \textbf{Yield strength $\sigma_y$}: stress at which irreversible deformation occurs.
\end{itemize}

In IBFE:
\begin{itemize}
    \item $E$ determines stiffness of structural elements and time-step limits.
    \item $\nu$ governs incompressibility; many biological tissues have $\nu \approx 0.45$–0.49.
    \item $\sigma_y$ is important when modeling non-linear or viscoplastic materials.
\end{itemize}


%%%%%%%%%%%%%%%%%%%%%%%%%%%%%%%%%%%%%%%%%%%%%%%%%%%%%%%
\subsection*{8. Compare linear elastic, nonlinear elastic, viscoelastic, and viscoplastic materials.}

\textbf{Answer.}


\begin{itemize}
    \item \textbf{Linear elastic}: $\sigma \propto \epsilon$; material returns to original shape.
    \item \textbf{Nonlinear elastic}: stress–strain relation is nonlinear but reversible.
    \item \textbf{Hyperelastic}: large deformations with a strain-energy function $W$.
    \item \textbf{Viscoelastic}: time-dependent response (creep, stress relaxation).
    \item \textbf{Viscoplastic}: time-dependent permanent deformation.
\end{itemize}

IBFE must choose the appropriate constitutive law depending on the biological system:
soft tissues often require nonlinear or viscoelastic models rather than linear elasticity.


%%%%%%%%%%%%%%%%%%%%%%%%%%%%%%%%%%%%%%%%%%%%%%%%%%%%%%%
\subsection*{9. What boundary conditions are commonly used in solid mechanics?}

\textbf{Answer.}

\begin{enumerate}
    \item Prescribed displacement
    \item Pressure or traction
    \item Mixed conditions
    \item Body forces (e.g., gravity)
    \item Contact between solids
    \item Thermal or process-induced stresses
\end{enumerate}

IBFE uses:
\begin{itemize}
    \item Lagrangian solid BCs to anchor or fix material points,
    \item Eulerian traction conditions for fluid–structure coupling,
    \item Contact or self-contact when structural meshes interact.
\end{itemize}


%%%%%%%%%%%%%%%%%%%%%%%%%%%%%%%%%%%%%%%%%%%%%%%%%%%%%%%
\subsection*{10. What is the importance of isotropy vs anisotropy in biological materials?}

\textbf{Answer.}

\begin{itemize}
    \item \textbf{Isotropic linear elastic}: material properties identical in all directions.
    \item \textbf{Anisotropic}: directional dependence (e.g., bone, tendon, myocardium).
\end{itemize}

In IBFE:
\begin{itemize}
    \item Isotropy simplifies constitutive laws.
    \item Anisotropy is required for fiber-reinforced tissues or materials whose stiffness differs along fiber directions.
\end{itemize}



%%%%%%%%%%%%%%%%%%%%%%%%%%%%%%%%%%%%%%%%%%%%%%%%%%%%%%%
\subsection*{11. How do stress and strain generalize in 3D, and why do we need tensors?}

\textbf{Answer.}


Thus:
\begin{itemize}
    \item Stress is a tensor describing forces acting on all possible surface orientations.
    \item Strain is a tensor describing how material line elements stretch, shear, or rotate.
\end{itemize}

Tensors allow:
\begin{itemize}
    \item compact representation of multidirectional relationships,
    \item coordinate-invariant properties,
    \item implementation in IBFE finite-element formulations where the deformation gradient and Cauchy stress appear naturally.
\end{itemize}

%%%%%%%%%%%%%%%%%%%%%%%%%%%%%%%%%%%%%%%%%%%%%%%%%%%%%%%
\end{document}